\Needspace{11\baselineskip}
\begin{center}
\textit{Die Anzahl der ambigen Klassen.}
\end{center}

\Article{257}
Unter allen Klassen in einer gegebenen Ordnung einer gegebenen Determinante verlangen besonders die ambigen Klassen eine ausführlichere Untersuchung; die Bestimmung ihrer Anzahl wird uns den Weg zu vielem anderen öffnen. Es genügt aber, diese Anzahl nur in der eigentlich primitiven Ordnung zu bestimmen, da sich die übrigen Fälle leicht auf diesen zurückführen lassen. Wir werden dazu zuerst zeigen, wie man alle eigentlich primitiven ambigen Formen \((A,B,C)\) der vorgelegten Determinante \(D\) findet, in denen entweder \(B=0\) oder \(B=\frac12 A\) gilt, und dann, wie man aus ihrer Anzahl die Anzahl aller eigentlich primitiven ambigen Klassen der Determinante \(D\) gewinnt.

\SourceLeaf{290}
I. Alle eigentlich primitiven Formen \((A,0,C)\) der Determinante \(D\) findet man offenbar, indem man für \(A\) die einzelnen Teiler von \(D\) (sowohl positiv als auch negativ) nimmt, für die \(C=-\frac{D}{A}\) zu \(A\) teilerfremd ist. Ist also \(D=1\), so gibt es zwei solche Formen, \((1,0,-1)\), \((-1,0,1)\); ebenso viele, wenn \(D=-1\), nämlich \((1,0,1)\), \((-1,0,-1)\); ist \(D\) eine Primzahl oder eine Primzahlpotenz (mit positivem oder negativem Vorzeichen), so gibt es vier: \((1,0,-D)\), \((-1,0,D)\), \((D,0,-1)\), \((-D,0,1)\). Ist allgemein \(D\) durch \(n\) verschiedene Primzahlen teilbar (zu denen hier auch \(2\) gerechnet werden muss), so gibt es insgesamt \(2^{n+1}\) solche Formen. Setzt man nämlich \(D=\pm PQR\ldots\), wobei \(P,Q,R\) usw. verschiedene Primzahlen oder Potenzen verschiedener Primzahlen bezeichnen, deren Anzahl \(=n\) ist, so sind die Werte von \(A\) gleich \(1,P,Q,R\) usw. und den Produkten aus beliebig vielen dieser Zahlen. Nach der Kombinationslehre ist die Anzahl dieser Werte \(2^n\); sie muss aber verdoppelt werden, da man jedem Wert sowohl ein positives als auch. ein negatives Vorzeichen geben muss.

II. Ebenso erhält man alle eigentlich primitiven Formen \((2B,B,C)\) der Determinante \(D\), indem man für \(B\) alle Teiler von \(D\) (positiv und negativ) nimmt, für die \(C=\frac12(B-\frac{D}{B})\) ganzzahlig und zu \(2B\) teilerfremd ist. Da \(C\) notwendig ungerade sein muss und daher \(CC\equiv1\ (\mathrm{mod.}\,8)\) gilt, folgt aus \(D=BB-2BC=(B-C)^2-CC\), dass \(D\) entweder \(\equiv3\ (\mathrm{mod.}\,4)\) ist, wenn \(B\) ungerade ist, oder \(\equiv0\ (\mathrm{mod.}\,8)\), wenn \(B\) gerade ist; ist also \(D\) einer der Zahlen \(1,2,4,5,6\) modulo \(8\) kongruent, so gibt es keine solchen Formen. Gilt \(D\equiv3\ (\mathrm{mod.}\,4)\), so ist \(C\) ganzzahlig und ungerade, welchen Teiler von \(D\) man auch für \(B\) nimmt; damit aber \(C\) mit \(2B\) keinen gemeinsamen Teiler hat, muss \(B\) so gewählt werden, dass \(\frac{D}{B}\) zu \(B\) teilerfremd ist. Daher gibt es für \(D=-1\) zwei Formen, \((2,1,1)\), \((-2.\ -1,-1)\), und allgemein erkennt man leicht: Ist die Anzahl aller Primzahlen, die \(D\) teilen, gleich \(n\), so entstehen insgesamt \(2^{n+1}\) Formen. --- Ist \(D\) durch \(8\) teilbar, so wird \(C\) ganzzahlig, wenn man für \(B\) irgendeinen geraden Teiler von \(\frac12 D\) nimmt; die andere Bedingung, dass \(C=\frac12 B-\frac{D}{2B}\) zu \(2B\) teilerfremd sein soll, erfüllt man \textit{erstens}, indem man für \(B\) alle Teiler von \(D\) nimmt, die das Doppelte einer ungeraden Zahl sind und für die \(\frac{D}{B}\) mit \(B\) keinen gemeinsamen Teiler hat. Ihre Anzahl (unter Berücksichtigung der verschiedenen Vorzeichen) ist \(2^{n+1}\), wenn \(D\) durch \(n\) verschiedene ungerade Primzahlen teilbar angenommen wird; \textit{zweitens}, indem man für \(B\) alle durch vier teilbaren Teiler von
\SourceLeaf{291}
\(\frac12 D\) nimmt, für die \(\frac{D}{2B}\) zu \(B\) teilerfremd ist. Ihre Anzahl ist ebenfalls \(2^{n+1}\), sodass es in diesem Fall insgesamt \(2^{n+2}\) solche Formen gibt. Setzt man nämlich \(D=\pm2^\mu PQR\ldots\), wobei \(\mu\) einen Exponenten größer als \(2\) bezeichnet und \(P,Q,R\) verschiedene ungerade Primzahlen oder Potenzen solcher Primzahlen bezeichnen, deren Anzahl \(n\) ist, so kann man \textit{sowohl} für \(\frac12 B\) \textit{als auch} für \(\frac{D}{2B}\) die Werte \(1,P,Q,R\) usw. und die Produkte aus beliebig vielen dieser Zahlen mit positivem und negativem Vorzeichen nehmen.

Aus allem folgt: Wird \(D\) durch \(n\) verschiedene ungerade Primzahlen teilbar angenommen (mit \(n=0\), wenn \(D=\pm1\) oder \(\pm2\) oder eine Zweierpotenz vorliegt), so ist die Anzahl aller eigentlich primitiven Formen \((A,B,C)\), in denen \(B\) entweder \(0\) oder \(\frac12 A\) ist, gleich \(2^{n+1}\), wenn \(D\) entweder \(\equiv1\) oder \(\equiv5\ (\mathrm{mod.}\,8)\) ist; gleich \(2^{n+2}\), wenn \(D\equiv2,3,4,6\text{ oder }7\ (\mathrm{mod.}\,8)\); und schließlich gleich \(2^{n+3}\), wenn \(D\equiv0\ (\mathrm{mod.}\,8)\). Vergleichen wir diese Anzahl mit dem, was wir in Artikel 231 für die Anzahl aller möglichen Charaktere primitiver Formen der Determinante \(D\) angegeben haben, so sehen wir, dass jene in allen Fällen genau doppelt so groß ist wie diese. Ist \(D\) negativ, so gibt es unter jenen Formen offenbar ebenso viele positive wie negative.
